\section{Síntesis y ampliación}

\begin{enumerate}
    \item La exploración es el desafío central del RL; de $\epsilon$-greedy a UCB/Thompson Sampling se forma un conjunto de herramientas cada vez más refinado
    \item En el RL profundo es necesario aproximar la novedad mediante modelos de densidad o errores de predicción de redes
    \item Los métodos basados en la curiosidad deben tener cuidado con el problema de la televisión ruidosa
    \item La metaexploración trata el propio «cómo explorar» como una capacidad que puede aprenderse
    \item En aplicaciones reales como la robótica y los LLM, una exploración eficaz determina directamente la eficiencia del aprendizaje
\end{enumerate}

\subsection{Aplicación ampliada: llevar la meta-exploration a la enseñanza de la computación}

Al final, la clase también mostró una dirección poco común pero muy sugerente: usar las estrategias de exploración del meta-RL para la calificación automática y la retroalimentación sobre los programas de los estudiantes. El sistema aprende primero «qué pruebas conviene ejecutar y qué señales de comportamiento aportan más información», y después encuentra los bugs y genera sugerencias de calificación con mayor rapidez.

\newpage
\begin{figure}[H]
\centering
\includegraphics[width=0.9\textwidth]{slides-images/slide-30.jpg}
\caption{Aplicación de la meta-exploration a la retroalimentación sobre programas: explorar qué interacciones exponen mejor los problemas del código de los estudiantes}
\end{figure}

\subsection{En qué se distingue la «exploración» en el contexto educativo del RL tradicional}

Aquí la exploración ya no consiste en controlar un robot o jugar un videojuego, sino en decidir: \textbf{para determinar dónde falla el programa de un estudiante, qué interacciones debe ejecutar primero el sistema, qué señales debe observar y qué fragmentos de video debe mostrar.} El valor de este tipo de exploración está en ahorrar tiempo a los asistentes de docencia y, al mismo tiempo, en hacer que la retroalimentación sea más focalizada y más accionable.

\begin{figure}[H]
\centering
\includegraphics[width=0.9\textwidth]{slides-images/slide-31.jpg}
\caption{Learned exploration en la tarea Bounce: interacciones distintas exponen categorías de errores distintas}
\end{figure}

\begin{table}[H]
\centering
\begin{tabular}{p{0.2\textwidth} p{0.28\textwidth} p{0.28\textwidth}}
\toprule
Tarea educativa & Información que hay que explorar & Beneficio para el sistema \\
\midrule
Retroalimentación sobre tareas de programación interactivas & Qué escenarios de prueba provocan bugs con mayor probabilidad & Localizar más rápido el origen de los errores \\
Apoyo a la calificación de los asistentes de docencia & Qué evidencias bastan mejor para respaldar la calificación según la rubric & Reducir la revisión manual y los juicios repetidos \\
Análisis de patrones de error & Qué comportamientos de los estudiantes pertenecen al mismo tipo de fallo & Formar plantillas de retroalimentación más reutilizables \\
\bottomrule
\end{tabular}
\caption{Objetivos de recolección de información de la meta-exploration en la enseñanza de CS}
\end{table}

\begin{figure}[H]
\centering
\includegraphics[width=0.9\textwidth]{slides-images/slide-32.jpg}
\caption{Práctica de AI-assisted grading en el curso: un flujo de calificación de los asistentes de docencia más rápido y más preciso}
\end{figure}

\subsection{Tabla general de esta clase}

\begin{table}[H]
\centering
\begin{tabular}{p{0.18\textwidth} p{0.22\textwidth} p{0.22\textwidth} p{0.22\textwidth}}
\toprule
Enfoque & Idea central & Ventajas & Problemas principales \\
\midrule
$\epsilon$-greedy / UCB / Thompson & Usar la incertidumbre para guiar la exploración en los bandits & Analizable, con regret demostrable & Difícil de extender directamente a MDP de gran escala \\
Deep RL intrinsic exploration & Fomentar la exploración con la novedad o el error de predicción & Se adapta a espacios de estados continuos y de alta dimensión & Fácil de engañar con ruido y con novedad espuria \\
End-to-end meta-exploration & Aprender la exploración y la ejecución directamente a partir del reward de la tarea & Principio unificado, óptimo en teoría & El problema de coupling dificulta el entrenamiento \\
DREAM-style decoupling & Aprender primero la información clave de la tarea y después la exploración que la recupera & Más fácil de optimizar, buenos resultados en la práctica & Depende del diseño de la representación de la tarea y del cuello de botella de información \\
\bottomrule
\end{tabular}
\caption{Mapa de los métodos de exploración de la Lecture 14}
\end{table}

\begin{knowledgebox}{Lecturas adicionales}
\begin{itemize}
    \item Auer et al., ``Finite-time Analysis of the Multiarmed Bandit Problem,'' Machine Learning 2002
    \item Bellemare et al., ``Unifying Count-Based Exploration and Intrinsic Motivation,'' NeurIPS 2016
    \item Burda et al., ``Exploration by Random Network Distillation,'' ICLR 2019
    \item Pathak et al., ``Curiosity-driven Exploration by Self-Supervised Prediction,'' ICML 2017
    \item Liu et al., ``Explore then Execute: Adapting without Rewards via Factorized Meta-Reinforcement Learning,'' 2022
\end{itemize}
\end{knowledgebox}

\end{document}
